% LayerTree talk: colours, beamer theme and the tree-drawing macros.
%
% The palette is the explainer site's (site/css/tokens.css): cream paper,
% brown-black ink, one vermilion signal colour for the proof path.

\definecolor{ltbg}{HTML}{F4EDE0}
\definecolor{ltsurface}{HTML}{FBF7EF}
\definecolor{ltsurface2}{HTML}{EDE4D3}
\definecolor{ltink}{HTML}{2A2119}
\definecolor{ltmuted}{HTML}{5E5043}
\definecolor{ltborder}{HTML}{D9CDB8}
\definecolor{ltstrong}{HTML}{8A7B69}
\definecolor{ltaccent}{HTML}{D1481A}
\definecolor{ltaccentsoft}{HTML}{F8DDD0}
\definecolor{ltnode}{HTML}{2767AE}
\definecolor{ltnodesum}{HTML}{1E4E85}
\definecolor{ltleaf}{HTML}{1D8A6A}
\definecolor{ltexit}{HTML}{A67300}
\definecolor{ltchain}{HTML}{6B4E9B}
\definecolor{ltgood}{HTML}{2E7D32}
\definecolor{ltgoodsoft}{HTML}{DFECD6}
\definecolor{ltbad}{HTML}{A3201B}
\definecolor{ltbadsoft}{HTML}{F6D9D3}
\definecolor{ltempty}{HTML}{BCAE98}

% ---------------------------------------------------------------- theme
\setbeamercolor{background canvas}{bg=ltbg}
\setbeamercolor{normal text}{fg=ltink,bg=ltbg}
\setbeamercolor{frametitle}{fg=ltink,bg=ltbg}
\setbeamercolor{title}{fg=ltink}
\setbeamercolor{subtitle}{fg=ltmuted}
\setbeamercolor{structure}{fg=ltaccent}
\setbeamercolor{alerted text}{fg=ltaccent}
\setbeamercolor{block title}{fg=ltink,bg=ltsurface2}
\setbeamercolor{block body}{fg=ltink,bg=ltsurface}
\setbeamercolor{item}{fg=ltaccent}
\setbeamercolor{footline}{fg=ltmuted}
\setbeamerfont{frametitle}{size=\Large,series=\bfseries}
\setbeamerfont{title}{size=\huge,series=\bfseries}
\setbeamertemplate{navigation symbols}{}
\setbeamertemplate{itemize items}{\color{ltaccent}\raisebox{0.15ex}{\scriptsize$\blacksquare$}}
\setbeamertemplate{itemize subitem}{\color{ltstrong}\raisebox{0.2ex}{\tiny$\blacksquare$}}
\setbeamertemplate{blocks}[rounded][shadow=false]
\setbeamertemplate{frametitle}{%
  \vspace*{0.35cm}%
  \begin{beamercolorbox}[wd=\paperwidth,leftskip=0.6cm]{frametitle}%
    \usebeamerfont{frametitle}\insertframetitle%
    \ifx\insertframesubtitle\@empty\else\\[-0.1ex]{\normalsize\color{ltmuted}\mdseries\insertframesubtitle}\fi%
  \end{beamercolorbox}%
  \vspace*{-0.15cm}%
  {\color{ltaccent}\rule{1.4cm}{1.4pt}}\hspace*{-1.4cm}\hspace*{0.6cm}%
}
\setbeamertemplate{footline}{%
  \hbox to \paperwidth{\hspace*{0.6cm}{\tiny\color{ltmuted}\texttt{LayerTree}\quad\insertsectionhead}\hfill{\tiny\color{ltmuted}\insertframenumber}\hspace*{0.5cm}}\vspace*{0.25cm}}
\setbeamertemplate{section page}{%
  \vfill
  \begin{center}
    {\color{ltaccent}\rule{2cm}{2pt}}\\[0.6em]
    {\Huge\bfseries\insertsectionhead}\\[0.4em]
  \end{center}
  \vfill}
\AtBeginSection{\begin{frame}[plain]\sectionpage\end{frame}}

\newcommand{\code}[1]{{\ttfamily\small #1}}
\newcommand{\codes}[1]{{\ttfamily\scriptsize #1}}
\newcommand{\tr}[1]{\textsf{\textbf{#1}}}
\newcommand{\muted}[1]{{\color{ltmuted}#1}}
\newcommand{\hl}[1]{{\color{ltaccent}\textbf{#1}}}
\newcommand{\lnote}[1]{\par\vspace{0.3em}{\small\color{ltmuted}#1}}

% ---------------------------------------------------------- overlay keys
% Style keys that change with the overlay: visible on=<2->, alt=<2>{a}{b}.
\tikzset{
  invisible/.style={opacity=0,text opacity=0},
  visible on/.style={alt={#1{}{invisible}}},
  alt/.code args={<#1>#2#3}{\alt<#1>{\pgfkeysalso{#2}}{\pgfkeysalso{#3}}},
}

% ---------------------------------------------------------- tree styles
\tikzset{
  lt/edge/.style={draw=ltstrong,line width=0.6pt},
  lt/pathedge/.style={draw=ltaccent,line width=1.6pt},
  lt/inode/.style={circle,draw=ltnode,fill=ltsurface,line width=0.8pt,
    minimum size=0.62cm,inner sep=0pt,font=\scriptsize\bfseries,text=ltnodesum},
  lt/ianc/.style={lt/inode,draw=ltaccent,line width=1.4pt},
  lt/ifold/.style={lt/inode,draw=ltaccent,fill=ltaccent,text=white,line width=1.4pt},
  lt/isib/.style={lt/inode,draw=ltaccent,dashed,line width=1.2pt,fill=ltaccentsoft},
  lt/iempty/.style={lt/inode,draw=ltempty,text=ltempty,dashed},
  lt/root/.style={rectangle,rounded corners=2pt,draw=ltnode,fill=ltsurface,
    line width=0.9pt,minimum height=0.6cm,minimum width=1.3cm,inner sep=2pt,
    font=\scriptsize\bfseries,text=ltnodesum},
  lt/rootanc/.style={lt/root,draw=ltaccent,line width=1.4pt},
  lt/rootfold/.style={lt/root,fill=ltaccent,text=white,draw=ltaccent},
  lt/leafbase/.style={rectangle,rounded corners=2pt,line width=0.9pt,
    minimum width=1.12cm,minimum height=0.78cm,inner sep=1pt,align=center,
    font=\scriptsize},
  lt/leaf-n/.style={lt/leafbase,draw=ltleaf,fill=ltsurface},
  lt/leaf-e/.style={lt/leafbase,draw=ltempty,dashed,fill=ltbg,text=ltempty},
  lt/leaf-new/.style={lt/leafbase,draw=ltleaf,fill=ltgoodsoft,line width=1.4pt},
  lt/leaf-chg/.style={lt/leafbase,draw=ltaccent,fill=ltaccentsoft,line width=1.4pt},
  lt/leaf-out/.style={lt/leafbase,draw=ltexit,fill=ltexit!15,line width=1.4pt},
  lt/leaf-op/.style={lt/leafbase,draw=ltstrong,fill=ltsurface2},
  lt/leaf-pend/.style={lt/leafbase,draw=ltchain,fill=ltchain!12,line width=1.2pt},
  lt/leaf-land/.style={lt/leafbase,draw=ltnode,fill=ltnode!10,line width=1.2pt},
  lt/leaf-chan/.style={lt/leafbase,draw=ltleaf,fill=ltleaf!15,line width=1.2pt,double,double distance=0.6pt},
  lt/leaf-sib/.style={lt/leafbase,draw=ltaccent,dashed,fill=ltaccentsoft,line width=1.2pt},
  lt/region/.style={rounded corners=5pt,fill=ltnode!8,draw=ltnode!50,dashed,line width=0.6pt},
  lt/badge/.style={rectangle,rounded corners=2pt,inner sep=1.5pt,
    font=\tiny\bfseries,text=white},
  lt/badge-sig/.style={lt/badge,fill=ltleaf},
  lt/badge-op/.style={lt/badge,fill=ltstrong},
  lt/badge-hash/.style={lt/badge,fill=ltchain},
  lt/badge-voucher/.style={lt/badge,fill=ltnode},
  lt/badge-warn/.style={lt/badge,fill=ltbad},
  lt/badge-exit/.style={lt/badge,fill=ltexit},
  lt/chain/.style={rectangle,rounded corners=3pt,fill=ltchain,text=white,
    font=\scriptsize\bfseries,inner sep=4pt,align=center},
  lt/chainpend/.style={rectangle,rounded corners=3pt,draw=ltchain,dashed,
    line width=1pt,fill=ltsurface,text=ltchain,font=\scriptsize\bfseries,inner sep=4pt,align=center},
  lt/tx/.style={rectangle,rounded corners=3pt,draw=ltink,line width=0.8pt,
    fill=ltsurface,font=\scriptsize,align=left,inner sep=5pt,text width=3.6cm},
}

% Geometry: leaf i of a depth-3 tree sits at (i*\LTdx, 0); the node of
% height h and index j at the centre of its leaves, h*\LTdy up.
\newcommand{\LTdx}{1.38}
\newcommand{\LTdy}{1.05}
\newcommand{\LTx}[2]{\fpeval{(#2*2^(#1) + (2^(#1)-1)/2)*\LTdx}}

% \LTtree[keys]{a/5/n/note, b/3/n/, -/0/e/, ...}  (exactly 8 entries)
%   label/amount/style/note; style is one of n e new chg out op pend land chan sib.
% keys: path=p (the proof of position p: ancestors and siblings),
%       path2=q (a second proof), fold (ancestors filled: recomputed),
%       region=h/i (an aligned subtree highlighted), root text, out text,
%       pending (output drawn dashed: computed, not committed),
%       sibleaf (draw the leaf sibling of the path as a proof element).
\pgfkeys{/ltk/.is family,/ltk,
  path/.store in=\LTkpath, path=-1,
  path2/.store in=\LTkpathB, path2=-1,
  fold/.store in=\LTkfold, fold=0,
  region/.store in=\LTkregion, region={},
  regionlabel/.store in=\LTkregionlabel, regionlabel={},
  rootname/.store in=\LTkrootname, rootname={r},
  out/.store in=\LTkout, out={},
  pending/.store in=\LTkpending, pending=0,
  nosums/.store in=\LTknosums, nosums=0,
}

\makeatletter
\newcommand{\LTsetamt}[2]{\expandafter\xdef\csname lt@amt@#1\endcsname{#2}}
\newcommand{\LTamt}[1]{\csname lt@amt@#1\endcsname}
% Is node (h,j) an ancestor of position p? Sets \LTtmp to 1 or 0.
\newcommand{\LTisanc}[3]{% h j p
  \ifnum#3<0 \def\LTtmp{0}\else
  \edef\LTtmp{\fpeval{ (floor(#3/2^(#1)) == #2) ? 1 : 0 }}\fi}
\newcommand{\LTissib}[3]{% h j p
  \ifnum#3<0 \def\LTtmp{0}\else
  \edef\LTtmp{\fpeval{ ((floor(#3/2^(#1)) - 2*floor(floor(#3/2^(#1))/2) == 0 ?
      floor(#3/2^(#1))+1 : floor(#3/2^(#1))-1) == #2) ? 1 : 0 }}\fi}

\newcommand{\LTtree}[2][]{%
  \pgfkeys{/ltk,path=-1,path2=-1,fold=0,region={},regionlabel={},rootname={r},out={},pending=0,nosums=0,#1}%
  % Parse the leaves.
  \def\lt@i{0}%
  \foreach \lab/\amt/\sty/\nt [count=\ii from 0] in {#2} {%
    \LTsetamt{\ii}{\amt}%
    \expandafter\xdef\csname lt@lab@\ii\endcsname{\lab}%
    \expandafter\xdef\csname lt@sty@\ii\endcsname{\sty}%
    \expandafter\xdef\csname lt@nt@\ii\endcsname{\nt}%
  }%
  % Region background.
  \ifx\LTkregion\@empty\else
    \expandafter\lt@region\LTkregion\@nil
  \fi
  % Edges.
  \foreach \h in {1,2,3} {%
    \pgfmathtruncatemacro{\nj}{2^(3-\h)-1}%
    \foreach \j in {0,...,\nj} {%
      \foreach \c in {0,1} {%
        \pgfmathtruncatemacro{\ch}{\h-1}%
        \pgfmathtruncatemacro{\cj}{2*\j+\c}%
        \LTisanc{\ch}{\cj}{\LTkpath}\let\ea\LTtmp
        \LTisanc{\ch}{\cj}{\LTkpathB}\let\eb\LTtmp
        \ifnum\numexpr\ea+\eb\relax>0
          \draw[lt/pathedge] (\LTx{\h}{\j},\h*\LTdy) -- (\LTx{\ch}{\cj},\ch*\LTdy);
        \else
          \draw[lt/edge] (\LTx{\h}{\j},\h*\LTdy) -- (\LTx{\ch}{\cj},\ch*\LTdy);
        \fi
      }%
    }%
  }%
  % Sums of interior nodes.
  \foreach \j in {0,1,2,3} {%
    \pgfmathtruncatemacro{\a}{\LTamt{\the\numexpr2*\j\relax}+\LTamt{\the\numexpr2*\j+1\relax}}%
    \expandafter\xdef\csname lt@s@1@\j\endcsname{\a}}%
  \foreach \j in {0,1} {%
    \pgfmathtruncatemacro{\a}{\csname lt@s@1@\the\numexpr2*\j\relax\endcsname+\csname lt@s@1@\the\numexpr2*\j+1\relax\endcsname}%
    \expandafter\xdef\csname lt@s@2@\j\endcsname{\a}}%
  \pgfmathtruncatemacro{\a}{\csname lt@s@2@0\endcsname+\csname lt@s@2@1\endcsname}%
  \expandafter\xdef\csname lt@s@3@0\endcsname{\a}%
  % Interior nodes.
  \foreach \h in {1,2} {%
    \pgfmathtruncatemacro{\nj}{2^(3-\h)-1}%
    \foreach \j in {0,...,\nj} {%
      \def\st{lt/inode}%
      \edef\sm{\csname lt@s@\h @\j\endcsname}%
      \ifnum\sm=0 \def\st{lt/iempty}\fi
      \LTissib{\h}{\j}{\LTkpath}\ifnum\LTtmp=1 \def\st{lt/isib}\fi
      \LTissib{\h}{\j}{\LTkpathB}\ifnum\LTtmp=1 \def\st{lt/isib}\fi
      \LTisanc{\h}{\j}{\LTkpath}\ifnum\LTtmp=1 \ifnum\LTkfold=1 \def\st{lt/ifold}\else\def\st{lt/ianc}\fi\fi
      \LTisanc{\h}{\j}{\LTkpathB}\ifnum\LTtmp=1 \ifnum\LTkfold=1 \def\st{lt/ifold}\else\def\st{lt/ianc}\fi\fi
      \node[\st] (n-\h-\j) at (\LTx{\h}{\j},\h*\LTdy) {\ifnum\LTknosums=0 \sm\fi};
    }%
  }%
  % Root.
  \def\rst{lt/root}%
  \ifnum\LTkpath>-1 \ifnum\LTkfold=1 \def\rst{lt/rootfold}\else\def\rst{lt/rootanc}\fi\fi
  \node[\rst] (root) at (\LTx{3}{0},3*\LTdy) {\LTkrootname\ \ \csname lt@s@3@0\endcsname};
  % Leaves.
  \foreach \i in {0,...,7} {%
    \edef\sty{\csname lt@sty@\i\endcsname}%
    \LTissib{0}{\i}{\LTkpath}\let\sa\LTtmp
    \LTissib{0}{\i}{\LTkpathB}\let\sb\LTtmp
    \edef\full{\sty}%
    \ifnum\numexpr\sa+\sb\relax>0 \ifnum\pdfstrcmp{\sty}{e}=0 \else\edef\full{sib}\fi\fi
    \edef\lab{\csname lt@lab@\i\endcsname}%
    \edef\amt{\csname lt@amt@\i\endcsname}%
    \ifnum\pdfstrcmp{\sty}{e}=0
      \node[lt/leaf-e] (l-\i) at (\i*\LTdx,0) {$\varnothing$};
    \else
      \node[lt/leaf-\full] (l-\i) at (\i*\LTdx,0) {\textbf{\lab}\\[-1pt]\amt};
    \fi
    \edef\nt{\csname lt@nt@\i\endcsname}%
    \ifx\nt\@empty\else
      \node[font=\tiny,text=ltmuted,anchor=north,inner sep=1pt] at (l-\i.south) {\nt};
    \fi
  }%
  % The tree output on the chain.
  \ifx\LTkout\@empty\else
    \ifnum\LTkpending=1
      \node[lt/chainpend,anchor=south] (out) at (\LTx{3}{0},3*\LTdy+0.55) {\LTkout};
      \draw[ltchain,dashed,line width=0.8pt] (out.south) -- (root.north);
    \else
      \node[lt/chain,anchor=south] (out) at (\LTx{3}{0},3*\LTdy+0.55) {\LTkout};
      \draw[ltchain,line width=0.8pt] (out.south) -- (root.north);
    \fi
  \fi
}
\def\lt@region#1/#2\@nil{%
  \pgfmathsetmacro{\rx}{#2*2^(#1)*\LTdx-0.68}%
  \pgfmathsetmacro{\rX}{((#2+1)*2^(#1)-1)*\LTdx+0.68}%
  \pgfmathsetmacro{\rY}{#1*\LTdy+0.42}%
  \fill[lt/region] (\rx,-0.62) rectangle (\rX,\rY);
  \ifx\LTkregionlabel\@empty\else
    \node[font=\tiny\bfseries,text=ltnode,anchor=north west,inner sep=2pt] at (\rx,\rY) {\LTkregionlabel};
  \fi
}
\makeatother

% \LTbadge{pos}{text}{kind} under a leaf (row 1); \LTbadgeB for row 2.
\newcommand{\LTbadge}[3]{\node[lt/badge-#3,anchor=north] at ($(l-#1.south)+(0,-0.3)$) {#2};}
\newcommand{\LTbadgeB}[3]{\node[lt/badge-#3,anchor=north] at ($(l-#1.south)+(0,-0.62)$) {#2};}
% A badge on an interior node or the root.
\newcommand{\LTnbadge}[3]{\node[lt/badge-#3,anchor=south] at ($(#1.north)+(0,0.05)$) {#2};}

% Placeholder for an opcode (CCV + hypothetical Merkle-sum opcode) figure not yet measured.
\newcommand{\msuTBD}[1]{{\color{ltmuted}opcode: not yet measured}}
